%%%%%%%%%%%%%%%%%%%%%%%%%%%%%%%%%%%%%%%%%%%%%%%%%%%%%%%%%%%%%%%%%%%%
%        Matematica dolce
%
% Copyright 2016-22 Daniele Zambelli
%
%------------------------------
% Matematica dolce per i licei, volume 1
%
% Compilation:
% pdflatex --shell-escape m_d_licei_21_1.tex
%------------------------------
%
% This work may be distributed and/or modified under the
% conditions of the LaTeX Project Public License, either version 1.3
% of this license or (at your option) any later version.
% The latest version of this license is in
%   http://www.latex-project.org/lppl.txt
% and version 1.3 or later is part of all distributions of LaTeX
% version 2005/12/01 or later.
%
% This work has the LPPL maintenance status `maintained'.
%
% The Current Maintainer of this work is
% Daniele Zambelli - daniele.zambelli@gmail.com
%
% This work consists of the files:
%  -  m_d_licei_21_1.tex (this file)
%  -  createpdf.py
%  -  readme.md
%  -  the part of code of the files under chap/, img/  and lbr/ directories
%%%%%%%%%%%%%%%%%%%%%%%%%%%%%%%%%%%%%%%%%%%%%%%%%%%%%%%%%%%%%%%%%%%%

%========================
% Definizione delle directory
%========================
% \newcommand{\basedir}{matematicadolce/}
\newcommand{\depdir}{deposito/}
\newcommand{\magdir}{\depdir magazzino/}
\newcommand{\matdir}{\depdir materiali/}
\newcommand{\moddir}{modifiche/}                    % MODIFICATO

%========================
% Classe del documento
%========================
\documentclass[10pt,a4paper,openright]{\magdir matsweetmem}
% Per la versione in scala di grigio commentare la precedente e
% decommentare la riga seguente
% \documentclass[10pt,a4paper,openright,gray]{\magdir matsweetmem}

%========================
% Variabili del progetto
%========================
\input{\magdir variabili}

%========================
% Variabili per questo volume
%========================
\newcommand{\editore}{Matematicamente.it}
\newcommand{\serie}{Matematica \(C^3\)}
\newcommand{\titolo}{Complementi di Matematica}     % MODIFICATO
\newcommand{\pdftitolo}{Matematica dolce}           % MODIFICATO
\newcommand{\docvers}{\texttt{8.0.1}}
\newcommand{\edizione}{2021}
\newcommand{\Edizione}{Edizione}
\newcommand{\tipo}{ (versione completa)}
\newcommand{\descr}{Testo per il primo biennio \protect\\
                    della Scuola Secondaria di II grado \\ \null
                    istituto professionale \\       % MODIFICATO
                    servizi socio-sanitari}         % MODIFICATO
\newcommand{\oggi}{19 giugno 2021}                  % MODIFICATO
\newcommand{\mese}{giugno}                          % MODIFICATO
\newcommand{\anno}{2021}
\newcommand{\mcisbn}{}

%========================
% Lettura preambolo
%========================
\input{\magdir packages}

%========================
% Lettura definizioni
%========================
\input{\magdir definizioni}
\input{\magdir definizioni_tikz}

%========================
% Con o senza linguaggio di programmazione
%========================
\newif\ifcoding
%\codingtrue     % commentare questa linea se NON si vuole il coding
\codingfalse    % commentare questa linea se si vogliono le parti di coding

%========================
% Documento
%========================
\begin{document}
\frontmatter
\pdfbookmark{frontespizio}{frontespizio}
\pagenumbering{gobble}
\begin{titlingpage}
 \frntspz
\end{titlingpage}

% Copyright (c) 2015 Daniele Zambelli - daniele.zambelli@gmail.com

\pagenumbering{roman}
\thispagestyle{empty}
\pdfbookmark{colophon}{colophon}
{\setlength{\parindent}{0em}\small{
\begin{center}
{\large{\serie – \titolo}}

Copyright {\textcopyright} {\anno} \editore
\end{center}

\begin{wrapfloat}{figure}{I}{0pt}
\includegraphics[width=0.2\columnwidth]{img/by-sa.png}
\end{wrapfloat}

Questo libro, eccetto dove diversamente specificato, è rilasciato nei 
termini della licenza Creative Commons Attribuzione – Condividi allo stesso 
modo 3.0 Italia (CC BY-SA 3.0) il cui testo integrale è disponibile agli 
indirizzi:\\
\url{http://creativecommons.org/licenses/by-sa/3.0/it/}\\
\url{http://creativecommons.org/licenses/by-sa/3.0/it/legalcode}

Tu sei libero:
di riprodurre, distribuire, comunicare al pubblico, esporre in pubblico, 
rappresentare, eseguire e recitare quest'opera, di modificare quest'opera, 
alle seguenti condizioni:

\emph{Attribuzione} --- Devi attribuire la paternità dell'opera nei modi 
indicati dall'autore o da chi ti ha dato l'opera in licenza e in modo tale 
da non suggerire che essi avallino te o il modo in cui tu usi l'opera.

\emph{Condividi allo stesso modo} --- Se alteri o trasformi quest'opera, 
o se la usi per crearne un'altra, puoi distribuire l'opera risultante solo 
con una licenza identica o equivalente a questa.

Per maggiori informazioni su questo particolare regime di diritto d'autore 
si legga il materiale informativo pubblicato in:\\
\url{http://www.copyleft-italia.it}.

\mcpar{Coordinatori del Progetto} \coord .

\mcpar{Autori} \autori.

\mcpar{Hanno Collaborato} \colab.

\mcpar{Progettazione e Implementazione in \LaTeX} {Dimitrios Vrettos}.

\mcpar{Collaboratori} {\texcol}.

\mcpar{Collaborazione, commenti e suggerimenti} Se vuoi contribuire anche tu alla stesura e aggiornamento 
del manuale Matematica Dolce o se  vuoi inviare i tuoi commenti e/o 
suggerimenti scrivi 
a~\mail{daniele.zambelli@gmail.com}.

\vspace{2ex}
 Versione del documento: {\docvers} del {\oggi}.

 Stampa edizione \edizione : \mese\ \anno.

 ISBN \mcisbn

\vspace{2ex}
 {\scshape{Dati tecnici per l'adozione del libro a scuola}}

 Titolo: \serie, \titolo\ -\edizione.

 Codice ISBN: \mcisbn 

 Editore: \href{http://www.matematicamente.it}{\editore}. 

 Anno di edizione: \anno.

 Prezzo pdf: \officialeuro\ 0,00.

 Formato: ebook (\scshape{pdf}).
}}
% \cleardoublepage

\intestazione{\matdir 00_intestazioni/les/}{indice}
\intestazione{\moddir 00_intestazioni/les/}{prefazione}
% ..................................................

%..................................................
%
% Secondo tema: Relazioni e Funzioni
%
\mainmatter
\parte{\moddir 0_a_parti/}{part_02_1_d}
\capitolo{\matdir 01_Logica/01_Insiemi_logica/ins01/}
         {insiemi}
\capitolo{\matdir 01_Logica/02_Relazioni/rel01/}
         {relazioni}
\capitolo{\matdir 07_Geometria_analitica/01_Punti_segmenti/pc01/}
         {pianocartesiano}
\capitolo{\matdir 07_Geometria_analitica/02_Retta/retta01/}
         {retta}
\capitolo{\matdir 07_Geometria_analitica/03_Parabola/para01/}
         {parabola_pc}
\capitolo{\matdir 04_Equazioni/06_Grado_2+/compl201/}
         {grasup2_irraz_valass}
% ..................................................
%
% Quarto tema: Dati e Previsioni
%
\parte{\moddir 0_a_parti/}{part_04_1_d}
\capitolo{\matdir 14_Statistica/01_Statistica_descrittiva/stat01/}
         {statistica}
\capitolo{\matdir 15_Probabilita/01_Calcolo_combinatorio/comb01/}
         {calc_combinatorio}
\capitolo{\matdir 15_Probabilita/02_Calcolo_probabilita/prob01/}
         {probabilita}
%..................................................
%
% Azzeramento numerazione capitoli
%
\renewcommand{\thechapter}{\Alph{chapter}}
\setcounter{chapter}{0}

\end{document}
